\documentclass[sigconf,anonymous]{acmart}
\settopmatter{printacmref=false}
\setcopyright{none}
\renewcommand\footnotetextcopyrightpermission[1]{}
\pagestyle{plain}

\usepackage{booktabs}
\usepackage{xcolor}
\usepackage{amsmath}
\newcommand{\pending}[1]{\textcolor{red}{[#1]}}
\newcommand{\ef}{\textit{ef}}
\newcommand{\KS}{\mathrm{KS}}

\usepackage{xr}
\externaldocument{main_anon}
\renewcommand{\thetable}{A\arabic{table}}
\renewcommand{\thefigure}{A\arabic{figure}}
\begin{document}
\title{Supplementary appendix to: PercEF: Exploiting Empirical Percentiles for Adaptive HNSW Search Beyond the Gaussian Assumption: [Experiments \& Analysis]}
\author{Anonymous Author(s)}
\affiliation{\institution{Anonymous}\country{}}
\maketitle
\appendix

\section{Additional results}

\begin{figure}[h]
  \centering
  \includegraphics[width=\linewidth]{figures/fig4_controlled.pdf}
  \caption{Ranking advantage of the empirical score against KS for the real datasets and for
  synthetic data with controlled anisotropy (B) and cluster separation (D).}
  \label{fig:controlled}
\end{figure}

\begin{figure}[h]
  \centering
  \includegraphics[width=\linewidth]{figures/fig3c_latency_cdf.pdf}
  \caption{Per-query latency on Deep1B and MS Turing (setting R). The empirical score spends less
  time on easy queries and more on hard ones; a fixed \ef{} spends about the same on all.}
  \label{fig:cdf}
\end{figure}

\begin{table*}[t]
  \caption{Time per distance computation, relative to fixed-\ef{} queries with the same distance count
  (setting P and R). Hard and easy: the 20\% of queries PercEF gives the largest and smallest \ef{},
  timed under the fixed \ef{} closest to PercEF's mean recall. Last column: PercEF's share of its
  distance computations spent on its hard queries.}
  \label{tab:querycost}
  \centering\footnotesize
  \resizebox{\linewidth}{!}{\input{tables/table3c_query_cost.tex}}
\end{table*}

\begin{table*}[h]
  \caption{Tail-weighted alternatives to KS on the raw vectors. Tail $z$: distance of the
  empirical upper quantile from the Normal's, in standard deviations. Tail mass: $\log_{10}$ of
  the observed over the expected share beyond the Normal's 99.9\% point. Margin:
  $|\rho_{\text{ours}}| - |\rho_{\text{Ada-ef}}|$, mean over P and R.}
  \label{tab:tail}
  \centering\footnotesize
  \input{tables/table5_tail_survey.tex}
\end{table*}

\begin{table*}[h]
  \caption{Offline cost, round-2 datasets.}
  \centering\footnotesize
  \input{tables/table4_offline.tex}
\end{table*}


\end{document}
